\documentclass[11pt,a4paper]{article}
\usepackage[margin=1in]{geometry}
\usepackage[utf8]{inputenc}
\usepackage{amsmath,amssymb}
\usepackage{booktabs}
\usepackage{tabularx}
\usepackage{microtype}
\usepackage{hyperref}
\usepackage{cite}

\title{\textbf{Adapting Language Models Under Severe GPU Memory Constraints:\\
Methods, Practical Trade-Offs, and Empirical Protocols}}
\author{Research Note}
\date{\today}

\begin{document}
\maketitle

\begin{abstract}
Fine-tuning large language models under strict GPU memory limits requires selecting adaptation methods whose trade-offs match available hardware. This note analyzes parameter-efficient adaptation, distinguishing architectural designs from empirical measurements across low-rank factorizations, bias tuning, dynamic rank allocation, quantization, and weight decomposition. We provide a comparative synthesis, isolate an unverified laboratory claim, and specify a controlled single-GPU experimental protocol.
\end{abstract}

\section{Memory Allocation and Parameter-Efficient Factorization}

During language model fine-tuning, GPU memory is consumed by weights, gradients, optimizer states, and activations. In 16-bit fine-tuning with Adam, each parameter requires 2 bytes for the weight, 2 bytes for the gradient, and 8 bytes for momentum buffers, demanding at least 16 bytes per parameter before allocating activations.

Parameter-efficient fine-tuning keeps pre-trained weights frozen. Low-Rank Adaptation (LoRA) parameterizes the linear update as $\Delta W = \frac{\alpha}{r} B A$ with rank $r \ll \min(d, k)$. LoRA freezes the pre-trained model weights and injects trainable rank decomposition matrices into each layer of the Transformer architecture, greatly reducing the number of trainable parameters for downstream tasks \cite{hu2021lora}.

Crucially, researchers must distinguish theoretical parameter savings from measured hardware footprints. Because frozen weights require no optimizer tracking, Adam memory drops to $O(r(d + k))$. Compared to GPT-3 175B fine-tuned with Adam, LoRA can reduce the number of trainable parameters by 10,000 times and the GPU memory requirement by 3 times \cite{hu2021lora}. However, this three-fold reduction was measured specifically on a 175B dense model trained with Adam; it cannot be transferred to smaller models or long-context regimens without direct measurement.

Localized bias tuning restricts updates further. In BitFit, most of the transformer-encoder parameters are frozen, and only the bias terms and the task-specific classification layer are trained \cite{benzaken2021bitfit}. In BERTBASE and BERTLARGE bias parameters make up 0.09\% and 0.08\% of the total number of parameters in each model, respectively \cite{benzaken2021bitfit}. Applying BitFit on pre-trained BERT models is competitive with fine-tuning the entire model \cite{benzaken2021bitfit}. While BitFit minimizes checkpoint storage, its capacity remains bounded on generation tasks.

\section{Dynamic Budgets, Quantization, and Weight Decomposition}

Standard LoRA assigns a uniform rank across layers. AdaLoRA parameterizes the incremental updates in the form of singular value decomposition \cite{zhang2023adalora}. Parameterizing updates as $\Delta W = P \Lambda Q$, AdaLoRA prunes redundant singular values via importance scoring while regularizing $P$ and $Q$. On the MNLI and SQuADv2.0 datasets, the performance of AdaLoRA under low budget levels can match the results of high budget settings \cite{zhang2023adalora}.

Even when optimizer states are minimized, 16-bit base weights create prohibitive memory footprints for larger models. QLoRA resolves this by compressing base weights into 4-bit NormalFloat (NF4), paired with Double Quantization and Paged Optimizers \cite{dettmers2023qlora}. Double Quantization saves an average of about 0.37 bits per parameter by quantizing the quantization constants \cite{dettmers2023qlora}. Operating across these components, QLORA reduces memory usage enough to finetune a 65B parameter model on a single 48GB GPU while preserving full 16-bit finetuning task performance \cite{dettmers2023qlora}.

Quantized base weights introduce deployment friction when merging adapters for serving. In QLoRA, post-training quantization (PTQ) of merged weights degrades accuracy. QA-LoRA introduces group-wise operators which increase the degree of freedom of low-bit quantization and decrease that of LoRA \cite{xu2023qa}. QA-LoRA consistently outperforms QLoRA with PTQ on top of LLMs of different scales and is on par with QLoRA without PTQ \cite{xu2023qa}.

Weight-Decomposed Low-Rank Adaptation (DoRA) investigates directional and magnitude dynamics. DoRA decomposes the pre-trained weight into two components, magnitude and direction, for fine-tuning, specifically employing LoRA for directional updates to efficiently minimize the number of trainable parameters \cite{liu2024dora}. By updating the directional and magnitude components of the QKV modules and only the magnitude of the rest of the layers, DoRA surpasses LoRA by 2.8\% on LLaMA-7B and 0.8\% on LLaMA-13B, while utilizing only less than half of the trainable parameters compared to LoRA \cite{liu2024dora}.

\section{Practical Trade-Offs and Comparative Synthesis}

Contrasting architectural designs with empirical findings clarifies key tradeoffs, summarized in Table~\ref{tab:comparison}.

\begin{table}[htbp]
\centering
\scriptsize
\caption{Systematic comparison of adaptation methods, explicitly categorizing architectural design specifications versus reported empirical benchmark results.}
\label{tab:comparison}
\begin{tabularx}{\textwidth}{l c c c X X}
\toprule
 & \multicolumn{3}{c}{\textbf{Architectural Design Specifications}} & \multicolumn{2}{c}{\textbf{Reported Empirical Benchmarks}} \\
\cmidrule(lr){2-4} \cmidrule(lr){5-6}
\textbf{Method} & \textbf{Precision} & \textbf{Tuned Scope} & \textbf{Zero Latency} & \textbf{Measured Memory Relief} & \textbf{Evaluated Model \& Benchmark Finding} \\
\midrule
Full FT & 16-bit & 100\% & Yes & $\ge 16$ B/param & Full parameter retraining baseline \\
BitFit \cite{benzaken2021bitfit} & 16-bit & $<0.1\%$ & Yes & Zero weight optimizer states & BERT: matches full FT on GLUE \\
LoRA \cite{hu2021lora} & 16-bit & 0.01--1\% & Yes & GPT-3: 1.2TB to 350GB & GPT-3: on par with full FT \\
AdaLoRA \cite{zhang2023adalora} & 16-bit & $\le 1\%$ & Yes & Pruned SVD parameterization & DeBERTaV3: matches full budget \\
QLoRA \cite{dettmers2023qlora} & 4-bit (NF4) & 0.1--2\% & No & 65B model on 48GB GPU & LLaMA-65B: matches 16-bit FT \\
DoRA \cite{liu2024dora} & 16-bit & 0.01--1\% & Yes & Magnitude-direction split & LLaMA-7B: $+2.8\%$ vs LoRA \\
QA-LoRA \cite{xu2023qa} & 4-bit (INT4) & 0.1--1\% & Yes & Group-wise 4-bit alignment & LLaMA: on par with QLoRA \\
\bottomrule
\end{tabularx}
\end{table}

A primary trade-off governs weight compression and throughput. QLoRA compresses static weights $4\times$, but dequantizing NF4 weights to BF16 per step reduces throughput. LoRA and DoRA preserve 16-bit throughput and allow zero-latency deployment by folding adapters directly into pre-trained weights.

\section{Proposed Single-GPU Empirical Protocol}

To evaluate these methods on consumer hardware, we specify a rigorous protocol. An informal lab assertion recently claimed: ``Our pilot halves training memory compared with LoRA.'' We treat this statement as explicitly unverified; lacking documented logs, it must not serve as an empirical result or design premise. Instead, it underscores the need for systematic peak VRAM profiling.

\subsection{Experimental Design and Controls}
\begin{itemize}
    \item \textbf{Hardware and Software Stack:} NVIDIA GeForce RTX 4090 24GB running Driver 550.54.14, CUDA 12.2, PyTorch 2.3.0, Transformers 4.40.0, PEFT 0.10.0, bitsandbytes 0.43.0, and FlashAttention-2 (v2.5.8).
    \item \textbf{Model and Data:} Meta-Llama-3-8B in bfloat16, adapted on the Alpaca dataset (52k examples) for 3 epochs (random seed 42).
    \item \textbf{Evaluated Methods:} 16-bit LoRA ($r=16, \alpha=32$), 4-bit QLoRA (NF4, double quantization), 16-bit DoRA ($r=16$), and 4-bit QA-LoRA ($L=32$).
    \item \textbf{Fixed Hyperparameters:} Per-device micro-batch size $b=1$ with gradient accumulation of 16 steps (effective batch 16); AdamW ($\beta_1=0.9, \beta_2=0.999$); learning rate $\eta = 2 \times 10^{-4}$ with linear schedule and 100 warmup steps.
    \item \textbf{Sequence Tiers:} Sequence lengths fixed at 512, 1024, and 2048 tokens; activation checkpointing and FlashAttention-2 enabled globally.
    \item \textbf{Instrumentation and Harness:} Peak VRAM tracked via PyTorch CUDA profiling (\texttt{max\_}\allowbreak\texttt{memory\_}\allowbreak\texttt{allocated}), steady-state throughput in tokens/second (steps 200--500), and downstream zero-shot accuracy on ARC-Challenge and HellaSwag using EleutherAI LM Evaluation Harness (v0.4.2).
\end{itemize}

\section{Limitations and Validity Boundaries}

Several empirical constraints govern this analysis. Peak GPU memory in parameter-efficient fine-tuning is often dominated by activation memory and sequence length rather than trainable parameter count alone. Intermediate activations outgrow optimizer buffers at long sequence lengths ($L \ge 2048$) unless activation checkpointing is applied. Reported benchmark gains on specific architectures such as LLaMA or BERT do not guarantee identical memory reductions or downstream accuracy on distinct model families without empirical verification. Finally, all six cited papers are registered and available in the project database; an automated registry warning indicated a title discrepancy between OpenAlex and Crossref for arXiv:2106.09685, which was verified as the canonical LoRA publication.

\bibliographystyle{plain}
\bibliography{../refs/references}

\end{document}
